% Generated draft: requires mathematical and source review.
\documentclass[11pt]{article}
\usepackage[margin=25mm]{geometry}
\usepackage{fontspec}
\setmainfont{Latin Modern Roman}
\usepackage{amsmath,amssymb,mathtools,unicode-math}
\setmathfont{Latin Modern Math}
\usepackage{enumitem,needspace,xurl,hyperref}
\setlength{\parindent}{0pt}
\setlength{\parskip}{5pt}
\setlength{\emergencystretch}{3em}
\newcommand{\N}{\mathbb{N}}
\newcommand{\Z}{\mathbb{Z}}
\newcommand{\Q}{\mathbb{Q}}
\newcommand{\R}{\mathbb{R}}
\newcommand{\C}{\mathbb{C}}
\newcommand{\F}{\mathbb{F}}
\newcommand{\A}{\mathbb{A}}
\newcommand{\PP}{\mathbb{P}}
\newcommand{\E}{\mathbb{E}}
\newcommand{\eps}{\varepsilon}
\newcommand{\dd}{\,\mathrm{d}}
\begin{document}
\section*{UnsolvedMath ID: 30004408}
\textbf{Source problem code:} catalogue entry
\subsection{Definitions and mathematical statement}
The supplied question is: ``As the number of tetrahedra tends to infinity among veering triangulations, does the proportion of non-layered veering triangulations tend to one?''

The record does not define the admissible veering triangulations, their equivalence relation for counting, or the predicate of being layered. Consequently, the following is a conditional counting formulation, pending specification of those choices.

Fix an admissible class \(\mathcal V\) of finite veering triangulations, an equivalence relation \(\sim\) preserving tetrahedron count, and a layeredness predicate invariant under \(\sim\). For a triangulation \(T\), let \(t(T)\) denote its number of tetrahedra. For each positive integer \(n\), define
\[
V_n=\{[T]_{\sim}:T\in\mathcal V,\ t(T)=n\},
\qquad
N_n=\{[T]_{\sim}\in V_n:T\text{ is non-layered}\}.
\]
Here non-layered means that the chosen layeredness predicate fails. Assume that every \(V_n\) is finite and that \(I=\{n\geq 1:V_n\neq\varnothing\}\) is unbounded. Under counting by exactly \(n\) tetrahedra, the question becomes whether
\[
\lim_{\substack{n\to\infty\\n\in I}}
\frac{|N_n|}{|V_n|}=1.
\]
Equivalently, for every \(\varepsilon>0\), does there exist a positive integer \(n_0\) such that
\[
|N_n|\geq(1-\varepsilon)|V_n|
\qquad\text{for every }n\in I\text{ with }n\geq n_0?
\]
The finiteness and unboundedness assumptions above are requirements for this formulation, not facts established by the supplied record.

The source also leaves unresolved whether counting is instead cumulative. With the same choices and assumptions, that interpretation asks whether
\[
\lim_{n\to\infty}
\frac{\sum_{k=1}^{n}|N_k|}{\sum_{k=1}^{n}|V_k|}=1,
\]
where the denominator is positive for all sufficiently large \(n\). The record does not select between these two counting interpretations. Neither formulation is fully specified until the admissible class, counting convention, and meaning of layeredness are fixed.
\subsection{Short English statement}
The question asks whether non-layered veering triangulations account for a proportion tending to one as tetrahedron count grows. A precise counting statement requires conventions for admissible triangulations, equivalence, layeredness, and exact-size versus cumulative counting; these are absent from the supplied record.
\subsection{Sources}
\begin{itemize}
\item Supplied local UnsolvedMath registry record: permanent ID \(30004408\), title ``Asymptotic Prevalence of Nonlayered Veering Triangulations,'' source problem code \texttt{catalogue entry}, in the Oberwolfach Reports Open Problems collection. The record gives proposed year \(2020\) and names no proposer. It labels the statement ``exact'' and describes it as clear, self-contained, publication-ready, and retained byte-for-byte; these are supplied editorial assessments, not independent verification of the missing counting conventions.
\item Supplied report citation: \emph{Low-dimensional Topology} (2021), Oberwolfach Reports, pp.\ 465--516, DOI \texttt{10.4171/owr/2020/8}. \url{https://doi.org/10.4171/owr/2020/8}. The background describes Oberwolfach Reports as documenting mathematical workshops and their open-problem discussions. The literature list calls this source ``Oberwolfach Report 8/2020.'' Both supplied date designations are retained.
\item Supplied literature reference: Futer, Gu\'{e}ritaud, and Schleimer, ``veering triangulations.'' \url{https://arxiv.org/abs/1911.00006}. This label and author attribution are reproduced as supplied; their correspondence to the linked item has not been independently checked.
\item Supplied literature review checked on \texttt{2026-08-22}: status ``open.'' Its assessment states that censuses and subsequent structural work give strong evidence and infinite non-layered families, but that no theorem establishes dominance of non-layered veering triangulations. This is supplied dated curation derived from the OWR material, not an independently verified literature assessment. The supplied background expressly requests independent verification of time-sensitive status and cited claims. The references listed for that assessment are the arXiv item and Oberwolfach report above.
\end{itemize}
\end{document}
